\label{sec:properties}

\subsection{Optimality: Minimising Relative Per-Capita Deviation}

The population pair test measures how far each pair of states is from
equal per-capita representation.
For states $i$ and $j$ with per-capita representations $r_i = p_i/s_i$
and $r_j = p_j/s_j$ (population per seat), the relative deviation is:
\[
  \delta(i,j) = \frac{|\,r_i - r_j|}{\min(r_i, r_j)}
              = \max\!\left(\frac{r_i}{r_j}, \frac{r_j}{r_i}\right) - 1.
\]

\begin{theorem}[Huntington-Hill Population Pair Test]
  \label{thm:hh-optimal}
  Among all apportionments satisfying $\sum s_i = H$ and $s_i \geq 1$,
  the Huntington-Hill apportionment minimises the maximum relative
  deviation $\max_{i \neq j} \delta(i, j)$.
\end{theorem}

\begin{proof}[Proof sketch]
  Suppose the HH apportionment is \emph{not} optimal and there exists
  a transfer of one seat from state $j$ to state $i$
  (changing $s_j \to s_j-1$, $s_i \to s_i+1$) that reduces
  $\max \delta$.
  For the transfer to reduce the deviation, it must be the case that
  giving state $i$ a seat improves equity more than it hurts
  state $j$.
  In HH, the seat allocation is determined by the priority
  $P_k(s_k) = p_k/\sqrt{s_k(s_k+1)}$; the next seat went to the
  state with highest priority.
  If state $j$ was awarded a seat over state $i$,
  then $P_j(s_j-1) > P_i(s_i)$, i.e.,
  $p_j/\sqrt{(s_j-1)s_j} > p_i/\sqrt{s_i(s_i+1)}$.
  One can verify that this priority inequality is equivalent to
  the condition that the pairwise deviation $\delta(i,j)$ in the
  HH allocation is no larger than in any allocation obtained
  by the transfer.
  A full proof by induction on the priority sequence appears in
  \citet[ch.~4]{balinski1982fair}.
\end{proof}

\citet{huntington1928method} originally proved this property, characterising
Huntington-Hill as the method that ``most nearly'' achieves equal
per-capita representation in the relative sense.
Webster's arithmetic mean criterion achieves the same property in
the absolute (not relative) sense.

\subsection{Paradox Immunity}

\begin{theorem}[Huntington-Hill Paradox Immunity]
  \label{thm:hh-paradox-free}
  The Huntington-Hill apportionment satisfies:
  \begin{enumerate}
    \item[(a)] \textbf{Alabama paradox immunity}: If $H$ increases to $H+1$
      (holding populations fixed), no state loses a seat.
    \item[(b)] \textbf{Population paradox immunity}: If state $i$'s
      population increases (holding $H$ and all other populations fixed),
      state $i$'s seat count does not decrease.
    \item[(c)] \textbf{New-states paradox immunity}: If a new state $k$
      is added with population $p_k$ and $\lfloor p_k/d \rfloor$ additional
      seats (where $d$ is the national divisor), no existing state's
      seat count changes.
  \end{enumerate}
\end{theorem}

\begin{proof}
  For (a): The HH priority-queue algorithm allocates seat $H+1$ to
  the state with the highest priority after seat $H$ has been allocated.
  This operation can only add a seat to one state; it never removes seats.
  No state's priority is re-evaluated to a lower value by the addition
  of a new seat.
  Hence the allocation for seats 1 through $H$ is identical whether
  we apportion $H$ or $H+1$ seats.

  For (b): If $p_i$ increases to $p_i' > p_i$, then every priority value
  $P_i(s) = p_i/\sqrt{s(s+1)}$ for state $i$ increases.
  This can only move state $i$ earlier in the priority sequence,
  meaning it receives its seats at least as early and thus at least
  as many seats.
  No other state's priorities change, so no other state's seat count
  increases; hence state $i$'s count cannot decrease
  (as the total is fixed at $H$, any gain by $i$ comes at the expense
  of another state, but not vice versa).

  For (c): If new state $k$ is added with $m = \lfloor p_k/d \rfloor$
  additional seats (i.e., $H$ increases to $H+m$ and $k$ gets exactly
  $m$ seats), the seats awarded to $k$ are exactly those $m$ steps in the
  priority queue that $k$ would occupy with population $p_k$.
  The divisor $d = N/(H) \approx (N+p_k)/(H+m)$ changes by $O(p_k/H^2)$,
  which for small $p_k/N$ does not perturb any other state's allocation.
  A formal proof requires showing that the priority queue steps 1 through
  $H+m$ include exactly the original $H$ steps for the $n$ existing
  states plus $m$ new steps for state $k$, interleaved by priority;
  see \citet[ch.~5]{balinski1982fair} for the full argument.
\end{proof}

\subsection{Bias Relative to Exact Proportionality}

Huntington-Hill is slightly small-state-favoring relative to Webster.
The geometric mean threshold $\sqrt{s(s+1)} > s + 0.5 - 1/(8s)$
(for all $s \geq 1$, with equality approached as $s\to\infty$)
implies that a state needs a slightly higher exact quota to ``earn''
a seat under HH than under Webster.
This counterintuitively means HH is slightly \emph{less}
favourable to large states (whose quotas are large integers) and
slightly \emph{more} favourable to small states (whose quotas are near
half-integers where the geometric vs.\ arithmetic mean difference matters).
The quantitative effect on 2020 Census data is examined in
Section~\ref{sec:comparison}.
